\section*{Resumen}
El presente Producto Integrador de Aprendizaje expone el diseño integral de un \textbf{Proyecto de Intervención Sustentable} formulado bajo la Metodología del Marco Lógico (MML) para mitigar la contaminación de mantos freáticos por escorrentía pecuaria y revertir la pérdida de carbono orgánico en suelos agrícolas del municipio de Marín y General Terán, Nuevo León. El proyecto articula la economía circular agropecuaria mediante la instalación de una Biofábrica Comunitaria de Co-compostaje termofílico y biodigestión anaeróbica para transformar las excretas bovinas y porcinas en biofertilizante sólido enriquecido y biol foliar. La intervención se alinea directamente con cuatro Objetivos de Desarrollo Sostenible de la Agenda 2030 (ODS 2 Hambre Cero, ODS 6 Agua Limpia, ODS 12 Producción Responsable y ODS 13 Acción por el Clima) y con las materias fundamentales de la Norma Internacional ISO 26000 sobre Gobernanza y Medio Ambiente. Mediante una matriz de indicadores objetivamente verificables y un cronograma operativo a 24 meses, se proyecta una reducción del $45\%$ en la carga de nitratos lixiviados, un incremento del $1.2\%$ en el contenido de materia orgánica edáfica y una disminución del $35\%$ en los costos de fertilización para 60 familias de productores ejidales.

\vspace{0.4cm}
\noindent\textbf{Palabras clave:} Responsabilidad social, desarrollo sustentable, biofertilización, economía circular, marco lógico, ODS 2030, Nuevo León.

\section{Introducción y Justificación Territorial}
La franja centro-noreste de Nuevo León combina una intensa actividad ganadera bovina y porcícola con valles agrícolas dedicados a granos básicos y citricultura \parencite{ocdeMexico2013, semarnat2020}. Sin embargo, la disposición inadecuada de estiércoles crudos en lechos de arroyos secos genera graves pasivos socioambientales: eutrofización de cuerpos de agua superficiales, emisión descontrolada de metano ($CH_4$) y óxido nitroso ($N_2O$), y proliferación de vectores de enfermedades zoonóticas \parencite{cantumartinez2020}.

Paralelamente, los suelos agrícolas de la región sufren una degradación acelerada, con niveles de materia orgánica frecuentemente inferiores al $1.0\%$, lo que merma su capacidad de retención de humedad en un contexto de sequía recurrente \parencite{ocdeMexico2024}. Frente a esta disyuntiva, la responsabilidad social universitaria exige formular proyectos de intervención que vinculen el conocimiento agronómico con las comunidades rurales para transformar un contaminante en un insumo regenerativo de alto valor agronómico \parencite{alfaro2019, onuIndicadores}.

\section{Diagnóstico Participativo y Árbol de Problemas}
Mediante talleres de diagnóstico rural rápido con productores de las comunidades de San José de González y El Castillo, se identificó el problema central: \textbf{Deterioro agroecológico y vulnerabilidad económica de pequeños productores por gestión lineal de residuos pecuarios e insumos caros}.
\begin{itemize}
    \item \textbf{Causas Raíz:} Ausencia de infraestructura de acopio, escasa capacitación técnica en bioprocesos de fermentación y cultura asistencialista desvinculada de la corresponsabilidad ambiental.
    \item \textbf{Efectos Negativos:} Pérdida de fertilidad del suelo, contaminación fecal de norias ejidales y dependencia crítica de fertilizantes nitrogenados importados de alto costo.
\end{itemize}

\section{Alineación con los Objetivos de Desarrollo Sostenible (ODS)}
El proyecto contribuye a metas específicas de la Agenda 2030 \parencite{onuIndicadores}:
\begin{enumerate}
    \item \textbf{ODS 2 (Meta 2.4):} Asegurar la sostenibilidad de los sistemas de producción de alimentos y aplicar prácticas agrícolas resilientes que mejoren la calidad de la tierra y el suelo.
    \item \textbf{ODS 6 (Meta 6.3):} Mejorar la calidad del agua reduciendo la contaminación, eliminando el vertimiento de excretas y minimizando la emisión de químicos peligrosos.
    \item \textbf{ODS 12 (Meta 12.5):} Reducir sustancialmente la generación de desechos mediante políticas de prevención, reducción, reciclado y reutilización biológica.
    \item \textbf{ODS 13 (Meta 13.3):} Mejorar la educación, sensibilización y capacidad comunitaria respecto a la mitigación del cambio climático y captura de carbono en el suelo.
\end{enumerate}

\section{Estructura Operativa: Biofábrica Comunitaria y Matriz de Marco Lógico}
El proyecto contempla el establecimiento de una Biofábrica Ejidal operada bajo esquema cooperativo:
\begin{table}[h!]
    \centering
    \small
    \begin{tabularx}{\textwidth}{lYY}
        \toprule
        \textbf{Nivel Jerárquico} & \textbf{Resumen Narrativo} & \textbf{Indicador Verificable (Metas)} \\
        \midrule
        \textbf{Fin} & Recuperar la salud de los agroecosistemas y la rentabilidad de las familias rurales en Nuevo León. & Incremento del $15\%$ en ingresos netos de las familias ejidales al año 3. \\
        \textbf{Propósito} & Valorizar agroecológicamente el $100\%$ de las excretas pecuarias en las comunidades piloto. & $3,600\text{ ton/año}$ de estiércol convertidas en $1,400\text{ ton}$ de compost grado A. \\
        \textbf{Componente 1} & Planta comunitaria de co-compostaje termofílico y biodigestión anaeróbica instalada. & Infraestructura civil y biofiltro al $100\%$ de capacidad operativa en mes 8. \\
        \textbf{Componente 2} & Programa de capacitación campesina en Buenas Prácticas Agrícolas y biofertilización. & 60 productores certificados y aplicando bioinsumos en 180 hectáreas. \\
        \textbf{Componente 3} & Monitoreo participativo de calidad de suelo, agua y balance de gases de efecto invernadero. & Cuatro muestreos anuales de $N, P, K$, MO edáfica y análisis microbiológico. \\
        \bottomrule
    \end{tabularx}
    \caption{Matriz sintética de Marco Lógico de la intervención agroecológica.}
\end{table}

\section{Viabilidad Técnica, Financiera y Responsabilidad Social (ISO 26000)}
\subsection{Proceso Biotecnológico}
El compostaje se gestiona en pilas aireadas de $1.5\text{ m}$ de altura con volteos mecánicos semanales. La fase termofílica ($55^\circ - 65^\circ\text{C}$) sostenida por más de 15 días garantiza la eliminación del $100\%$ de patógenos entéricos (\textit{Salmonella spp.}, \textit{Escherichia coli}) y semillas de malezas. La adición de rastrojo de maíz molido ajusta la relación carbono-nitrógeno ($C/N$) a una proporción óptima de $30:1$.

\subsection{Gobernanza y Sostenibilidad Comunitaria}
Siguiendo los lineamientos de la ISO 26000, el proyecto no es asistencialista: se constituye un Comité de Vigilancia Ciudadana con paridad de género para administrar el fondo revolvente de comercialización del excedente de biofertilizantes, asegurando la autonomía financiera y operativa del proyecto una vez concluido el acompañamiento universitario.

\section{Conclusiones}
1. La intervención sustentable demostró que el residuo pecuario constituye un activo estratégico para la soberanía alimentaria y la regeneración de suelos en Nuevo León.
2. La articulación de la Metodología de Marco Lógico con los ODS y la ISO 26000 confiere viabilidad técnica y transparencia social al proyecto.
3. Se consolida el compromiso de la Universidad Autónoma de Nuevo León y su Facultad de Agronomía de promover una ciencia agrícola humanista, socialmente responsable y ecológicamente comprometida.

